\section{Choose a display cell without changing the model}
After reflecting the raw position, select one of the eight display cells. The
exercise compares floor rounding with rounding to the nearest integer, resolving
an exact half upward. For a nonnegative raw position \(p\), these are
\[
  \operatorname{floorCell}(p)=\left\lfloor\frac p{256}\right\rfloor,
  \qquad
  \operatorname{nearestCell}(p)=\left\lfloor\frac{p+128}{256}\right\rfloor.
\]
The second rule is specifically for the renderer's nonnegative domain. Applying
the same adjustment to arbitrary negative values would define a different signed
tie rule and needs its own explanation.

Adding 128 adds half a pixel before dropping the fractional bits. At raw 127,
the nearest cell is zero. At raw 128, it is one. At raw 129, it is also one.
The upper coordinate raw 1792 still selects cell seven under either rule. The
renderer cannot produce cell eight from a valid position.

\Needspace{14\baselineskip}
\tagpdfsetup{table/header-rows={1}}
\begin{longtable}{rrrr}
  \toprule
  Raw position & Position in pixels & Floor cell & Nearest cell \\
  \midrule
  127          & \(127/256\)        & 0          & 0            \\
  128          & \(1/2\)            & 0          & 1            \\
  192          & \(3/4\)            & 0          & 1            \\
  320          & \(1+1/4\)          & 1          & 1            \\
  448          & \(1+3/4\)          & 1          & 2            \\
  1664         & \(6+1/2\)          & 6          & 7            \\
  1792         & 7                  & 7          & 7            \\
  \bottomrule
\end{longtable}
The initial raw position \((320,448)\) becomes cell \((1,1)\) with floor and
\((1,2)\) with nearest rounding. Those are pixel IDs nine and seventeen in an
eight-column row-major frame. The raw position is the same in both runs.

The renderer receives the raw values and a rounding-mode argument. It writes one
white cell in an otherwise black 64-word frame and returns that cell's ID. Its
cell coordinates are derived values. Do not write them back as the next raw
positions. Doing so discards the accumulated fractions and changes the motion.

Nearest rounding can display the ball in wall cell seven while its raw coordinate
is only 6.5. That displayed cell does not mean a collision has occurred. A later
step may still have a positive velocity even while the screen cell remains seven.
The collision rule uses the raw interval, and the renderer observes its result.
Bouncing from the rounded cell would give the two display modes different physical
trajectories.
